\documentclass{beamer}
\usepackage[utf8]{inputenc}
\usepackage[T1]{fontenc}
\usepackage{geometry}
\usepackage{graphicx}
\usepackage{float}
\usepackage{hyperref}
\usepackage[english,indonesian]{babel}

\title{Optimasi Pemrosesan Template LaTeX Menggunakan Pendekatan Berbasis AI}
\author{Academic Support Agent}
\date{\today}

\begin{document}
\maketitle

\begin{abstract}

\end{abstract}

\chapter{PENDAHULUAN}


\chapter{TINJAUAN PUSTAKA}


% Original content (commented out):
% % IMPORTANT! In order for the document to compile, one needs to use XeLaTeX or LuaLaTeX as compiler. This can be done in  Overleaf by Menu -> Settings -> Compiler -> Choose XeLaTeX/LuaLaTeX
\documentclass[t,24pt]{beamer}

\usepackage{KUstyle}
\usepackage{multicol} % Package for multiple coloumns

\toplinje{Text at top} % The text at top. Remove the command if no text is desired

\begin{document}

% The first slide. One can for instance change the main title, the subtitle, speaker, KU-unit, date and backgroundimage
{
\setbeamertemplate{background}{\includegraphics[width=\paperwidth,height=\paperheight]{pictures/frontpage.jpg}}
\begin{frame}
    \begin{textblock*}{\textwidth}(0.57\textwidth,0.1\textheight)
        \begin{beamercolorbox}[wd=6.4cm,ht=7.7cm,sep=0.5cm]{hvidbox}
            \fontsize{4}{10}\fontfamily{ptm}\selectfont \textls[200]{UNIVERSITY OF COPENHAGEN}
            \noindent\textcolor{KUrod}{\rule{5.4cm}{0.4pt}}
        \end{beamercolorbox}
    \end{textblock*}
    \begin{textblock*}{\textwidth}(0.57\textwidth,0.1\textheight)
        \begin{beamercolorbox}[wd=6.4cm,sep=0.5cm]{hvidbox}
                \Huge \textcolor{KUrod}{Main title}
                \vspace{0.5cm}
                \par
                \Large This is the subtitle
                \vspace{0.5cm}
                \par
                \normalsize Speaker, KU-unit, date
        \end{beamercolorbox}
    \end{textblock*}
    \begin{textblock}{1}(14.2,11.44)
        \includegraphics[width=1cm]{KU/KU-logo.png}
    \end{textblock}
\end{frame}
}

% A standard slide. It's important that either [hvid], [rod] or [billede] is include af \begin{frame}. [hvid] makes a white background, [rod] makes a red background and [billede] adds an image to the background. It you use [billede] it's important to change the \setbeamertemplate{background} command (see below how one does so)
\begin{frame}[hvid]
\frametitle{White standard slide}
Lorem ipsum dolor sit amet, consectetur adipiscing elit. Quisque viverra rhoncus urna, et sodales magna efficitur et. Integer tempor augue eget ultricies malesuada. Praesent maximus sapien et tellus venenatis, nec consequat sem dictum. In lectus tellus, sagittis at massa in, tempor pretium quam. Pellentesque vel scelerisque mi.
\\~\\
Sed vel ligula egestas, commodo lacus in, cursus velit. Nam faucibus felis et nibh vehicula, eget aliquam justo facilisis. Aenean a dui sed libero malesuada dictum sit amet quis turpis. Nulla consectetur lorem at augue lacinia, ut rutrum neque vulputate. Aliquam ultrices suscipit felis vel pellentesque. Donec lacinia orci nec feugiat rhoncus. Nunc congue at velit vitae posuere. In hac habitasse platea dictumst.
\end{frame}

% This is examples of different types of slides
\begin{frame}[hvid]
\frametitle{List with bullets}
    \begin{itemize}
        \item Bullet 1
        \item Bullet 2
            \begin{itemize}
                \item Bullet 2.1
                \item Bullet 2.2
            \end{itemize}
        \item Bullet 3
    \end{itemize}
\end{frame}

\begin{frame}[hvid]
\frametitle{Slide with 2 coloumns}
    \begin{multicols}{2}
        \begin{itemize}
            \item Bullet 1
            \item Bullet 2
        \end{itemize}
        \vfill\null
        \columnbreak
        \begin{itemize}
            \item Bullet 3
            \item Bullet 4
            \item Bullet 5
        \end{itemize}
    \end{multicols}
\end{frame}

\begin{frame}[hvid]
\frametitle{Slide with coloumn and image}
    \begin{multicols}{2}
        \begin{itemize}
            \item Bullet 1
            \item Bullet 2
        \end{itemize}
        \vfill\null
        \columnbreak
            \includegraphics[width=5.1cm,height=6.1cm,keepaspectratio]{pictures/1.jpg}
    \end{multicols}
\end{frame}

\begin{frame}[hvid]
\frametitle{Slide with large image}
    \vspace{0.5cm}
    \includegraphics[width=11.3cm,height=6.4cm,keepaspectratio]{pictures/2.jpg}
\end{frame}

% This is an example with [billede], where one adds a backgroundimage
{
\setbeamertemplate{background}{\includegraphics[width=\paperwidth,height=\paperheight]{pictures/3.jpg}}
\begin{frame}[billede]
    \begin{textblock*}{\textwidth}(0\textwidth,0.77\textheight)
        \begin{beamercolorbox}[wd=6.4cm,sep=0.3cm]{rodbox}
        Some text
        \end{beamercolorbox}
    \end{textblock*}
\end{frame}
}

{
\setbeamertemplate{background}{\includegraphics[width=\paperwidth,height=\paperheight]{pictures/4.jpg}}
\begin{frame}[billede]
    \begin{textblock*}{\textwidth}(0.57\textwidth,0.4\textheight)
        \begin{beamercolorbox}[wd=6.4cm,sep=0.3cm]{rodbox}
        Here is some text and a list.
            \begin{itemize}
                \item Bullet 1
                \item Bullet 2
                \item Bullet 3
                \item Bullet 4
            \end{itemize}
        \end{beamercolorbox}
    \end{textblock*}
\end{frame}
}

\begin{frame}[rod]
    \fontsize{80}{0}\selectfont ``
    \\
    \vspace{-0.6cm}
    \fontsize{32}{0}\selectfont Prediction is very difficult, especially if it's about the future.
    \\
    \vspace{2.8cm}
    \fontsize{18}{0}\selectfont Niels Bohr
\end{frame}

\begin{frame}[rod]
    \vspace{0.9cm}
    \resizebox{\textwidth}{!}{
        \begin{tabular}{r l}
            \textbf{Research} & Forskning \\
            \textbf{Education} & Uddannelse \\
            \textbf{Exchange of knowdledge} & Forskningsformidling \\
        \end{tabular}
    }
\end{frame}

\end{document}

\end{document}